\section{Step 4: Response}\label{sec:response}

\emph{What a per-flow FPR hides.} A per-flow Benign FPR of 3.8\% looks modest in a table; at a gateway that blocks host pairs, it means blocked innocent hosts. Garcia and Valeros note that a few false positives can have strong consequences once detections drive IP blocking and that earlier alerts help only at a controlled FPR (Sec.~6.3--6.4)~\cite{garcia2023labeling}, and Akbari et al.\ that, because detectors are prone to false positives, simply dropping suspected traffic can have severe consequences~\cite{akbari2020atmos}. We therefore score response in false blocks, counted per \emph{host-pair window}, a host pair within one 60\,s test window. A block is false if the blocked pair carries no attack flow in that window; the pairs seen in a window without an attack flow are the eligible benign host-pair windows.

\emph{ARP watcher.} For the attack that flow features cannot see (Section~\ref{sec:splits}), a small watcher raises a high alert when a MAC other than the known one announces the router IP (.1), the classic arpwatch idea~\cite{arpwatch}. In the MITM-ARP window it raised one high alert 1.7\,s into the window, and none in the other nine windows. This shows that the rule works here, not a false-alarm rate.

\emph{Policies.} With enforcement on, the gateway blocks a host pair (OVS drop rules for both directions) once the pair has collected $k$ attack verdicts; the harness resets these counters after each window. Optionally, it ignores MITM-ARP flow verdicts, which the ARP watcher covers. Requiring $k$ verdicts is our design choice. We meant to select the policy on validation data, as the one with the fewest false blocks among the ten policies ($k\in\{1,2,3,5,10\}$, MITM-ARP verdicts ignored or not) within 0.005 of the best share of attack flows that start after their pair is blocked. The rule was degenerate: all ten policies lie within 0.0012 of the best share, so it reduced to the fewest false blocks and picked the largest $k$ in the grid ($k=10$, MITM-ARP ignored; 1 false block on validation, against 128 for $k=1$). The result is therefore the trade-off over $k$ (Table~\ref{tab:kcurve}); the enforced policy is one point on it.

\begin{table}[t]
\centering
\caption{Trade-off over $k$, simulated per 60\,s window on the verdicts logged in run~1. Ign.: MITM-ARP verdicts ignored. False: false blocks; /h: per hour of capture time. Class columns: attacker pairs blocked, of the number in parentheses (BF BruteForce, DS DDoS-SYN, DH DoS-HTTP, MA MITM-ARP, Mi Mirai, Re Recon); Att.: their sum. Time: median time from a blocked attacker pair's first flow to its block.}
\label{tab:kcurve}
\footnotesize
\setlength{\tabcolsep}{2.4pt}
\begin{tabular}{rcrrcccccccrr}
\hline
 & & & & BF & DS & DH & MA & Mi & Re & Web & Att. & Time \\
$k$ & Ign. & False & /h & (1) & (6) & (1) & (7) & (7) & (1) & (3) & (26) & (s) \\
\hline
1  & no  & 64 & 384 & 1 & 6 & 1 & 0 & 7 & 1 & 3 & 19 & 19.0 \\
1  & yes & 37 & 222 & 1 & 6 & 1 & 0 & 7 & 1 & 3 & 19 & 19.0 \\
2  & yes &  5 &  30 & 1 & 6 & 1 & 0 & 6 & 1 & 3 & 18 & 33.8 \\
3  & yes &  1 &   6 & 1 & 6 & 1 & 0 & 1 & 1 & 3 & 13 & 15.8 \\
5  & yes &  0 &   0 & 1 & 6 & 1 & 0 & 0 & 1 & 1 & 10 & 15.3 \\
10 & yes &  0 &   0 & 1 & 6 & 1 & 0 & 0 & 1 & 0 &  9 & 15.3 \\
20 & yes &  0 &   0 & 0 & 6 & 1 & 0 & 0 & 1 & 0 &  8 & 15.3 \\
\hline
\end{tabular}
\end{table}

\emph{The trade-off.} Table~\ref{tab:kcurve} is simulated on a run without enforcement, so its $k=1$ counts differ slightly from the enforced run; per-hour rates extrapolate from ten minutes of traffic. Without ignoring MITM-ARP verdicts, false blocks fall from 384 per hour at $k=1$ to 84, 12 and 0 at $k=2$, 3 and 5; ignoring them alone cuts the $k=1$ count from 64 to 37. The DDoS-SYN, DoS-HTTP and Recon attackers are blocked at every $k$, so $k$ trades false blocks against the low-rate attackers. Most Mirai bots are lost at $k=3$ (1 of 7 blocked) and all from $k=5$; the web attackers drop to 1 of 3 at $k=5$ and to none from $k=10$; the brute-force attacker is lost at $k=20$. No MITM-ARP pair is blocked at any $k$. These losses hold within 60\,s windows with per-window counters, where each run saw only 4--5 live Mirai flows; with counters kept over a longer horizon, both attackers and benign pairs would reach a given $k$ more often.

\begin{table}[t]
\centering
\caption{Enforced live blocking over the ten test windows (final harness, blocks cleared per window). Ign.: MITM-ARP verdicts ignored. Attackers: attacker pairs blocked. False: false blocks of eligible benign host-pair windows, and per hour of capture time. Dropped: packets.}
\label{tab:blocking}
\footnotesize
\setlength{\tabcolsep}{3pt}
\begin{tabular}{lrrrrr}
\hline
Policy & Blocked & Attackers & False & /h & Dropped \\
\hline
$k=1$ & 79 & 19/26 & 60/1,674 & 360 & 326,268 \\
$k=10$, ign. & 9 & 9/26 & 0/1,675 & 0 & 163,546 \\
\hline
\end{tabular}
\end{table}

\emph{Enforced runs.} Blocking on the first verdict (Table~\ref{tab:blocking}) blocked all 7 Mirai bots (about 100k packets, 35\% of the window, dropped) but also 60 of 1,674 eligible benign host-pair windows (3.6\%), i.e.\ 360 false blocks per hour of capture time, involving 90 distinct IPs, in 9 of the 10 windows. Most fell in the DDoS-SYN (12.5\%) and DoS-HTTP (6.8\%) windows; 21 were triggered by MITM-ARP verdicts, 14 of them outside the MITM-ARP window. Each cuts a benign pair off for the rest of the window. The $k=10$ policy blocked all six DDoS attackers (133k packets dropped), the DoS-HTTP attacker (30k), the brute-force attacker (217) and the port scanner (0; the scan had ended), with no false blocks, and missed the 7 Mirai bots and the 3 web attackers. In the $k=1$ run the DDoS-SYN window was replayed at only 52\% of the original rate; since slower replay splits flows (Section~\ref{sec:fidelity}), its flow and block counts in that window may be inflated.

\emph{Late blocks.} In the run-1 simulation only 11.1\% of the attack flows start after their pair is blocked, at every $k$: a block needs a verdict, and a verdict needs a flow to end. Blocks still cut the later packets of long attack flows: the enforced runs dropped 0.503 ($k=1$) and 0.342 ($k=10$) of the DDoS-SYN window's packets and 0.707 and 0.698 of the DoS-HTTP window's. The validation share (0.995) is thus far too optimistic, as validation ignores the export delay.
